\documentclass[12pt]{article}

\usepackage{amsmath}
\usepackage{parskip}
\usepackage{amssymb}
\usepackage{geometry}
\usepackage{hyperref}
\usepackage{booktabs}
\usepackage{hyperref}
\usepackage{tabularx}
\usepackage{float}
\usepackage{graphicx}
\usepackage{longtable}
\usepackage{url}
\usepackage{natbib}
\usepackage{pdfpages}


\geometry{a4paper, margin=1in}
\title {Nichols Radial Injection Model (RIM) IX\\
\newline
{Extreme Hawking Decoupling: Observational Evidence from Sgr A*}}
\author{Lawrence William Nichols}
\date{\today}

\begin{document}

\maketitle
\begin{abstract}
This paper presents the physical foundation for a proposed ``Phase Separation'' mechanism operating at the event horizon of the supermassive black hole Sgr A*. We examine the co-existence of heavy, sub-relativistic baryonic material---including atomic species (H, He, C, N, O, Ne, Mg, Si, Fe) and molecular complexes (HCN, CO, H$_2$O)---with highly ionized, high-energy electromagnetic emission in the Galactic Center. Within the Nichols Radial Injection Model (RIM), these heavy components constitute a ``Slow Mass Phase'' that, upon approaching relativistic regimes near the horizon, undergoes  decoupling from a ``Fast Photon Phase'' associated with Fe XXV/XXVI X-ray flares. This process is proposed to generate ``widowed'' baryonic mass, interpreted here as a candidate for dark matter, while the decoupled electromagnetic component is trapped in the black hole.
\end{abstract}

\section{The Sgr A* Laboratory: Baryonic Inventory}

The immediate environment of Sgr A* offers a unique observational window into material interacting with a supermassive black hole. Multi-wavelength data, particularly from the Herschel Space Observatory and the Atacama Large Millimeter/sub-millimeter Array (ALMA), together with X-ray spectroscopy, reveal the simultaneous presence of both massive, sub-relativistic components and high-energy electromagnetic emission.

\subsection{The Mini-Spiral: Pre-Decoupling Matter}

The mini-spiral structure in the Galactic Center serves as the primary reservoir feeding material toward Sgr A*. Observations have established the presence of a rich baryonic inventory, including:

\begin{itemize}
    \item \textbf{Atomic species:} Hydrogen (H), helium (He), carbon (C), nitrogen (N), oxygen (O), neon (Ne), magnesium (Mg), silicon (Si), and iron (Fe).
    \item \textbf{Molecular complexes:} Hydrogen cyanide (HCN), carbon monoxide (CO), and water (H$_2$O).
\end{itemize}

In the RIM framework, these molecular and atomic components represent baryonic matter in a coupled, pre-decoupling state. As this material spirals inward under gravitational influence, tidal forces and increasing thermal energies are expected to dissociate molecular complexes into their constituent atomic nuclei, priming them for the proposed separation process at the horizon.

\subsection{The Flare Events: Release of the Photon Phase}

In contrast to the sub-relativistic bulk material in the mini-spiral, X-ray observations frequently detect energetic flares characterized by highly ionized iron lines, notably Fe XXV and Fe XXVI. These features are interpreted within the RIM as signatures of the ``Fast Photon Phase'' released during decoupling.

\begin{itemize}
    \item The electromagnetic component (massless radiation) propagates at the speed of light $c$ and cannot sustain stable orbits, leading to rapid infall or ejection along null geodesics.
    \item The associated baryonic nuclei, possessing rest mass ($m > 0$), remain constrained to sub-luminal velocities, creating a kinematic divergence that drives  separation.
\end{itemize}

\subsection{Phase Separation at Relativistic Thresholds}

Ordinary baryonic matter remains coupled in its atomic and molecular form while spiraling inward at sub-relativistic speeds in the mini-spiral. However, as material is accelerated toward the event horizon in the Radial Injection Model, it approaches velocities of order $v \approx 0.9c$. At such relativistic regimes, a fundamental kinematic constraint emerges: massive particles cannot reach the speed of light $c$, while the associated electromagnetic radiation propagates exactly at $c$.

This velocity threshold enforces a ``phase separation'' analogous to Stephen Hawking's heuristic description of black hole radiation. Hawking explained that quantum mechanics predicts virtual particle-antiparticle pairs constantly materializing and annihilating throughout space. Near a black hole's event horizon, gravitational effects can separate the pair: one member may fall inward while the other escapes to infinity, appearing as radiation emitted by the black hole. In this picture, the horizon induces a divergence in the fates of the paired entities.

Extending this separation principle classically to real baryonic material within the RIM, we propose the following process:

As infalling matter reaches $v \approx 0.9c$, the Lorentz factor $\gamma \approx 2.29$ produces significant time dilation. The heavy, massive nuclei (Fe, Si, Mg, etc.) remain sub-luminal and cannot maintain the null geodesics required for stable photon-like orbits. Consequently, the system undergoes a ``velocity divorce'':

\begin{itemize}
    \item The \textbf{light/fast component} (electromagnetic radiation associated with highly ionized flares, such as Fe XXV and Fe XXVI) loses any effective rest mass. It can no longer sustain an orbit and is forced along null geodesics into the event horizon or into the black hole.
    \item The \textbf{heavy/slow component} (baryonic nuclei with $m > 0$) cannot attain $c$. This kinematic mismatch causes  decoupling: the massive fraction is ejected outward in relativistic outflows or jets, now stripped of its electromagnetic signature (``EM-less'').
\end{itemize}

The ejected ``widowed'' baryonic mass, traveling at high velocity and subject to extreme time dilation, falls out of temporal synchronization with the observable universe. Its gravitational influence persists, yet it remains effectively invisible electromagnetically, providing a candidate explanation for dark matter phenomenology.

This mechanism implies that Sgr A* functions not merely as an accretor but as an active ``Extreme Hawking Decoupling Engine,'' sorting matter according to its relativistic behavior rather than consuming it whole.

This velocity disparity is central to the proposed ``velocity divorce'' mechanism.

\section{ Decoupling: The Velocity Divergence}

Within the RIM, the region near the event horizon functions analogously to a centrifugal separator, where orbital dynamics and relativistic effects enforce separation between massive and massless components. Because stable orbits require nonzero rest mass, the pathways of matter and radiation diverge fundamentally.

\subsection{Lorentz Factor and Time Dilation Effects}

Heavy elements such as Fe, Si, and Mg observed or inferred in outflows reach velocities approaching $v \approx 0.9c$. The corresponding Lorentz factor is given by
\begin{equation}
\gamma = \frac{1}{\sqrt{1 - \frac{v^2}{c^2}}} \approx 2.29.
\end{equation}

This relativistic time dilation is proposed to desynchronize the ejected ``widowed'' baryonic material from the local cosmic time frame, rendering its gravitational influence effectively invisible on conventional timescales and thereby contributing to the observed effects attributed to dark matter.

\section{Conclusion}

The concurrent detection of complex molecular species and atomic material in the mini-spiral alongside highly ionized iron flares from Sgr A* provides compelling observational context for a phase-separation engine operating at the black hole horizon. In the Radial Injection Model, Sgr A* does not simply accrete or consume infalling matter; instead, it  sorts baryonic components according to their relativistic behavior. The massive fraction is ejected into the galactic environment as ``widowed'' material, while the electromagnetic component is channeled either into the black hole and lost or escapes and is ejected into space.

This framework offers a novel interpretation of existing multi-wavelength data and suggests new directions for testing phase-separation signatures through coordinated observations across radio, infrared, and X-ray regimes.

\section{Video Catalog}
A video overview of the full RIM series is available on the author's \href{https://www.youtube.com/@Nichols_Thought_Experiments}{Nichols Radial Injection Model (RIM)} YouTube channel.

\begin{thebibliography}{99}

\bibitem{rimI}
L.\,W.\,Nichols.
Nichols Radial Injection Model (RIM) and Radial ERB Inflows: A Mechanism for Progenitor-less Astrophysical Events, the Hubble Pulse, and 4D Substrate Interactions.
[Self-published], January 21, 2026.
\newline
\url{https://rxiverse.org/abs/2602.0036}

\bibitem{rimII}
L.\,W.\,Nichols.
Nichols Radial Injection Model (RIM) 11-Year Solar Cycle Pulses, GRB Milestones, Historical Alignments, and the 16.6 Gyr Hypersphere Universe.
[Self-published], February 2026.
\newline
\url{https://rxiverse.org/abs/2602.0038}

\bibitem{rimIII}
L.\,W.\,Nichols.
Nichols Radial Injection Model (RIM) III: The 3,000-Year Universal Odometer, the 1054 CE Primary Strike, and the Decaying 11-Year Cycle.
[Self-published], 15 February 2026.
\newline
\url{https://rxiverse.org/abs/2602.0043}

\bibitem{rimIV}
L.\,W.\,Nichols.
Nichols Radial Injection Model (RIM) IV: A 5D Mass-Regulated Manifold Solution to Cosmic Expansion.
[Self-published], 28 February 2026.
\newline
\url{https://rxiverse.org/abs/2603.0001}

\bibitem{rimV}
L.\,W.\,Nichols.
Nichols Radial Injection Model (RIM) V:
The Empirical Seal of the 11.07-Year Pulse
and Universal Resonance in Stellar Cycles
[Self-published], 4 March 2026.
\newline
\url{https://rxiverse.org/abs/2603.0003}

\bibitem{rimVI}
L.\,W.\,Nichols.
Nichols Radial Injection Model (RIM) VI:
The Birth of the Universe
[Self-published], 10 March 2026.
\newline
\url{https://rxiverse.org/abs/2603.0023}

\bibitem{rimVII}
L.\,W.\,Nichols.
Nichols Radial Injection Model (RIM) VII:
the ΛCDM’s Inconsistencies
[Self-published], 23 March 2026.
\newline
\url{https://rxiverse.org/abs/2603.0085}

\bibitem{rimVIII}
L.\,W.\,Nichols.
Nichols Radial Injection Model (RIM) VIII:
VIII The Bullet Cluster and Geometric Torsion
[Self-published], 29 March 2026.
\newline
\url{https://rxiverse.org/abs/2603.0102}

\bibitem{einstein1935}
Einstein, A., \& Rosen, N. (1935).
\textit{The Particle Problem in the General Theory of Relativity}.
Physical Review, 48(1), 73.

\end{thebibliography}

\newpage

\section{Engineering Studies and Visual Data}
This section includes relevant visual diagrams.
\includepdf[pages=-]{Universe_Mass.pdf}
\end{document}